\chapter{Theory}
What follows is going to be the introduction and explanation of specific concepts required to understand the key points of this thesis.
\section{Cavity Quantum Electrodynamics}
At the heart of an atom-photon gate lies the interaction between the two components, which in the context of this thesis is achieved by placing an atom in a high-finesse Fabry-Perot cavity.
The resulting dynamics then fall in the domain of Cavity Quantum Electrodynamics (CQED).
Placing the atom between two highly-reflective mirrors of an optical cavity increases the amount of times a photon that is inside of a cavity bounces back and forth.
This effectively increases the absorption cross section by the number of round-trips the photon takes, allowing for a strong light-matter interaction.

The theoretical model which explains this interaction in a fully quantum mechanical picture is called the \textit{Jaynes-Cummings Model} \cite{JaynesCummings1963} and was named after Edwin Jaynes and Fred Cummings who developed it in 1963.
\subsection{Jaynes-Cummings Model}

\begin{itemize}
	\item Basic setup is atom with 3 levels $\ket{u}$, $\ket{c}$, $\ket{e}$ and cavity with single mode that is resonant to $\ket{c} \leftrightarrow \ket{e}$, with detuning between resonant transition and cavity, include picture from \cite{RempeReiserer2015} and cite ofc
	\item Write down hamiltonian that includes first the basic structure and then explain rotating wave approximation (possibly show derivation in appendix)
	\item Show that the solutions of the eigenstates have this splitting term leading to the jaynes cummings ladder where at zero detuning (aka resonance)
	      levels are seperated by 2g
\end{itemize}

\subsection{Dissipation and Decay Channels}

Now we can include some losses here and there which are important to take into account as they reduce fidelity

Will prolly scrap this section and rather talk about the cavities and their birefringence

Should prolly just include these two subsubsections in the one before to save space and just include extra stuff that I'm concidering like mode matching, extra dump state, and whatnot

\subsubsection{Atomic Decay Channels}
\subsubsection{Photonic Decay Channels}
\section{Reflection from a Cavity-Atom System}

\begin{itemize}
	\item Different cases for reflection and thus transmission when driving with weak coherent pulse, based on state of the atom and detunings and stuff like that
	\item system was introduced by gardiner and collet \cite{GardinerColletIO1985} show equation and then say how you can solve for analytical solution for photonic field reflection $r(\omega)$ and then show the equation for the coupled and uncoupled case and maybe already show theoretical plots that compare the two both in reflectivity and phase shift.
\end{itemize}

\section{Conditional Reflection}

\begin{itemize}
	\item Go off what Kimble and Duan were saying in \cite{DuanKimble2004} and how you get a phase shift which can also be seen in the plot when plotting the phase shift across different detunings and explain where it's coming from
	\item Say that this can be used to build a controlled phase flip gate when photon is in eigenbasis of the cavity with the atom being the controlling part and the photon being the target part
\end{itemize}
\subsection{Reflection in CNOT basis}
\begin{itemize}
	\item Show that by turning the basis and instead of being in the polarization eigenbasis of the cavity, one can build a CNOT gate which is much more interesting for quantum computation by just sending a superposition of the two eigenbases of the cavity.
\end{itemize}
